\section{Related literature}\label{sec:lit}

\paragraph{Do low-frequency liquidity measures measure liquidity?} A long line of work benchmarks
spread and price-impact proxies built from daily data against estimates from intraday trades and
quotes. \citet{goyenko2009} establish the approach for U.S.\ equities; \citet{corwin2012} and
\citet{abdi2017} propose the high-low and close-high-low estimators that now dominate applied work.
\citet{jahanparvar2023} show that several of these estimators carry a volatility-induced bias that
dominates when the true spread is small relative to volatility, and that the bias has grown as
spreads have narrowed. In cryptocurrency markets, \citet{brauneis2021} compare low-frequency
measures with order-book benchmarks for bitcoin and ether on three exchanges and find that the
Corwin-Schultz and Abdi-Ranaldo estimators track time variation in liquidity best, while Amihud
and Kyle-Obizhaeva measures capture levels. We add to this literature in three ways. First, we score
proxies of \emph{signed order flow}, not only of spreads, against an exchange-reported ground truth
for twenty coins over almost five years. Second, we use a fee change as a controlled shock to the
composition of volume, which tests whether volume-based liquidity measures track liquidity or merely
volume. Third, we ask whether the measurement error matters for the downstream use that motivates
many applications, namely short-horizon return forecasting.

\paragraph{Trade classification and flow toxicity.} When aggressor flags are unavailable, trades are
signed by the tick rule or by \citet{leeready1991}. \citet{easley2012,easley2016} propose bulk volume
classification (BVC), which assigns a fraction $\Phi(\Delta p/\sigma)$ of each bar's volume to buyers,
and argue that it captures the information content of order flow better than trade-level rules.
Evidence on its accuracy is mixed. \citet{chakrabarty2015} find that BVC is less accurate than the
tick rule and Lee-Ready for individual equities, \citet{panayides2019} find that it identifies
aggressors less accurately but is more closely related to informed trading and improves after
calibration, and \citet{andersen2014} question whether VPIN built on BVC predicts toxicity events.
These studies use equity or futures data over comparatively short samples. We exploit the fact
that Binance reports the taker-buy volume of every bar to evaluate BVC on hourly and one-minute bars
over about 42{,}000 hours for each of twenty coins, and to bound how much of the true flow any
return-based classifier can recover.

\paragraph{Trading fees, fee-free volume and wash trading.} Theory predicts that trading fees shape
quotes, depth and the mix of traders \citep{colliard2012,foucault2013}, and natural experiments in
equity markets show that changes in fees and transaction taxes move spreads, depth, volume and market
composition \citep{malinova2015,colliard2017}. In cryptocurrency markets, \citet{brauneis2022} study
the introduction of maker fees on Coinbase Pro and find wider quoted spreads, less depth and fewer
but larger trades. Two recent studies use the Binance promotion that we also exploit.
\citet{galati2024} shows with order-book data that removing maker and taker fees led market makers to
widen spreads and reduce depth, so that total transaction costs rose despite the waiver, and raises
concerns about trading conduct; \citet{bozzetto2026} use the return of fees in March 2023 and
conclude that zero fees made order flow less informative about prices. We use the same event for a
different purpose, as a controlled shock to the composition of volume that reveals what bar-based
measures respond to, with nineteen untreated coins as controls and the restoration of fees as a
second test. The episode also connects to the literature on fake volume in cryptocurrency markets,
which detects wash trading with statistical tests \citep{cong2023,amiram2025} and with leaked
account-level records \citep{aloosh2024}. Aggressor flags do not identify the parties to a trade, so
we cannot label individual trades as wash trades. We show instead that the volume created by zero
fees was two-sided within the minute and within the second, and that bar-based measures read it as
liquidity.

\paragraph{Order flow and short-horizon returns.} Order imbalance predicts short-horizon index
returns \citep{chordia2002}, order-flow imbalance at the top of the book explains contemporaneous
price changes \citep{cont2014}, and deep-learning models extract predictive signals from
order-book flow at horizons of seconds to minutes \citep{kolm2023}. For cryptocurrencies,
\citet{makarov2020} document that signed order flow explains a large share of bitcoin returns across
exchanges, \citet{aleti2021} compare spot and futures microstructure, and
\citet{easley2024crypto}, following \citet{easley2021machine}, report that microstructure measures
such as the Roll measure and VPIN predict future crypto market dynamics in a machine-learning
framework. Time-series predictability at horizons of days to weeks is documented by
\citet{liu2021}. Our study differs from this work in focusing on the intraday horizon at which
microstructure effects should be strongest, in measuring predictive gains strictly out of sample
with tests designed for nested and overlapping forecasts \citep{clark2007,campbell2008}, and in
pricing the resulting signals at realistic transaction costs. Because a single dataset invites
specification search, we fix the design before estimation, report every feature set and model we
estimate, and interpret Sharpe ratios with the caution recommended by \citet{bailey2014} and
\citet{harvey2016}.
